\documentclass[11pt]{article}

\usepackage[margin=1in]{geometry}
\usepackage{amsmath,amssymb}
\usepackage{graphicx}
\usepackage{booktabs}
\usepackage{hyperref}
\usepackage{float}
\usepackage{caption}
\usepackage{subcaption}

\graphicspath{{figures/}}

\title{Residue-Class Tiling as a Finite Analogue for Sparse Feature Recovery}
\author{Dan Hawkley}
\date{}

\begin{document}

\maketitle

\begin{abstract}
We study feature recovery on a controlled finite system defined by Mod30 residue classes.
Using a continuous embedding of integers constrained to residues coprime with 30, we compare
non-negative matrix factorization (NMF) and top-$k$ sparse autoencoders (SAEs). We introduce
metrics for structural recovery, redundancy, and capacity utilization. While NMF achieves higher
coverage and efficient use of components, SAEs produce sparse and interpretable features but leave
substantial latent capacity unused or redundant. These results show that sparsity-induced
interpretability does not guarantee efficient representation, even with known ground truth.
\end{abstract}

\section{Introduction}

Understanding how models recover structure is difficult when ground truth is unknown. We construct
a finite system where structure is explicit:

\begin{itemize}
    \item integers mapped under modulo 30,
    \item filtered to residues coprime with 30,
    \item yielding 8 persistent residue classes.
\end{itemize}

This provides a controlled testbed where
\[
    \text{true features} = \text{residue lanes}.
\]

\section{Mod30 Residue Structure}

Let
\[
    R = \{1,7,11,13,17,19,23,29\}.
\]
These define the persistent Mod30 lanes. We embed filtered integers into a continuous space and
evaluate whether factorization and sparse-learning methods recover these lanes.

\section{Continuous Embedding}

Each integer is mapped to:
\begin{itemize}
    \item circular coordinates,
    \item lane channels with local overlap,
    \item additive noise.
\end{itemize}

\begin{figure}[H]
\centering
\includegraphics[width=0.55\textwidth]{35_shared_mod30_embedding.png}
\caption{Continuous embedding of Mod30 residue classes.}
\label{fig:shared-embedding}
\end{figure}

\section{Models}

\subsection{Non-negative Matrix Factorization}

NMF factorizes the continuous embedding matrix as
\[
    X \approx W H,
\]
where $W$ contains component activations and $H$ contains learned components.

\subsection{Sparse Autoencoder}

The top-$k$ sparse autoencoder is trained as
\[
    X \longrightarrow Z \longrightarrow \operatorname{top}\text{-}k(Z) \longrightarrow \widehat{X}.
\]
Only $k$ latent units remain active per sample after the top-$k$ mask.

\section{Metrics}

We evaluate both model families using the same recovery metrics:

\begin{itemize}
    \item \textbf{Coverage}: recovered lanes divided by total lanes.
    \item \textbf{Purity}: alignment of learned features with residue classes.
    \item \textbf{Reconstruction error}: relative error between $X$ and $\widehat{X}$.
    \item \textbf{Redundancy}: total activations divided by captured lanes.
    \item \textbf{Effective features}: captured lanes divided by total components.
    \item \textbf{Feature status}: dead, useful, and redundant feature fractions.
\end{itemize}

\section{Results}

\subsection{Coverage vs. Reconstruction}

\begin{figure}[H]
\centering
\includegraphics[width=0.72\textwidth]{36_nmf_vs_sae_coverage_vs_reconstruction.png}
\caption{NMF vs. SAE: coverage and reconstruction tradeoff.}
\label{fig:coverage-reconstruction}
\end{figure}

NMF achieves stronger reconstruction and higher coverage in this controlled benchmark.

\subsection{Sparsity vs. Coverage}

\begin{figure}[H]
\centering
\includegraphics[width=0.72\textwidth]{37_nmf_vs_sae_sparsity_vs_coverage.png}
\caption{Sparse activations vs. structural recovery.}
\label{fig:sparsity-coverage}
\end{figure}

The SAE enforces sparse activations, but sparsity alone does not guarantee complete recovery.

\subsection{Capacity Utilization}

\begin{figure}[H]
\centering
\includegraphics[width=0.82\textwidth]{38_nmf_vs_sae_effective_dead_redundant.png}
\caption{Effective features and dead or redundant capacity.}
\label{fig:capacity}
\end{figure}

SAEs leave significant latent units unused or redundantly allocated.

\subsection{Redundancy Tradeoff}

\begin{figure}[H]
\centering
\includegraphics[width=0.72\textwidth]{39_nmf_vs_sae_redundancy_vs_coverage.png}
\caption{Redundancy vs. coverage.}
\label{fig:redundancy}
\end{figure}

Redundancy separates efficient recovery from overlapping or diffuse feature allocation.

\subsection{Feature Alignment}

\begin{figure}[H]
\centering
\begin{subfigure}{0.46\textwidth}
    \centering
    \includegraphics[width=\textwidth]{40_nmf_alignment_8_components.png}
    \caption{NMF, 8 components}
\end{subfigure}
\hfill
\begin{subfigure}{0.46\textwidth}
    \centering
    \includegraphics[width=\textwidth]{41_sae_alignment_L8_k1.png}
    \caption{SAE, $L=8$, $k=1$}
\end{subfigure}
\caption{Feature-to-lane alignment for matched-capacity models.}
\label{fig:alignment}
\end{figure}

\subsection{Activation Patterns}

\begin{figure}[H]
\centering
\begin{subfigure}{0.46\textwidth}
    \centering
    \includegraphics[width=\textwidth]{42_nmf_activation_matrix_8.png}
    \caption{NMF}
\end{subfigure}
\hfill
\begin{subfigure}{0.46\textwidth}
    \centering
    \includegraphics[width=\textwidth]{43_sae_activation_matrix_L8_k1.png}
    \caption{SAE}
\end{subfigure}
\caption{Activation matrices showing dense factorization vs. sparse latent activity.}
\label{fig:activations}
\end{figure}

\section{Main Result}

The core tradeoff is:
\[
    \text{sparsity} \Longrightarrow \text{interpretability}
    \quad \text{but not necessarily} \quad
    \text{efficient capacity}.
\]

More precisely, dense factorization efficiently uses its components, while sparse autoencoders can
leave significant latent capacity unused or redundant even in overcomplete regimes.

\section{Discussion}

The experiments isolate a finite version of the efficiency--interpretability tradeoff:
\[
    \text{efficiency} \longleftrightarrow \text{interpretability}.
\]

Sparse models can improve clarity while reducing completeness. In this Mod30 benchmark,
interpretability must therefore be evaluated jointly with coverage, redundancy, and effective
feature utilization.

\section{Conclusion}

We introduced a finite Mod30 system for evaluating feature recovery and showed:

\begin{itemize}
    \item NMF recovers structure efficiently.
    \item SAEs produce interpretable but incomplete representations.
    \item Sparsity does not guarantee effective utilization.
\end{itemize}

\section{Reproducibility}

All results are generated via the following notebooks:

\begin{verbatim}
01_mod30_residue_manifold
02_local_tiling_vs_global_capture
03_sparse_feature_analogy
04_mod30_sae_toy_model
05_true_sparse_autoencoder
06_nmf_vs_sae_recovery_comparison
\end{verbatim}

\end{document}
